\documentclass{tex_aux/rithesis}
\usepackage[utf8]{inputenc}
\usepackage[T1]{fontenc}
\usepackage{geometry}
\usepackage{graphicx}
\usepackage{float}
\usepackage{hyperref}
\usepackage[english,indonesian]{babel}

\title{Optimasi Pemrosesan Template LaTeX Menggunakan Pendekatan Berbasis AI}
\author{Academic Support Agent}
\date{\today}

\begin{document}
\maketitle

\begin{abstract}

\end{abstract}

\chapter{PENDAHULUAN}


\chapter{TINJAUAN PUSTAKA}


% Original content (commented out):
% %% osvjezeno: 25.6.2024. (Miroslav Joler)

\documentclass[botnum,a4paper,11,oneside,final]{tex_aux/rithesis}
\usepackage[utf8]{inputenc}
\usepackage[T1]{fontenc}
\usepackage[croatian]{babel}
\input{tex_aux/rithesis_preamble}
%%%%%%%%%%%%%%%%%%%%%%%%%%%%%%%%%%%%%%%%%%%%%


% pomocu \includeonly moze se kompajlirati samo neko specificno poglavlje, da se skrati vrijeme kompajliranja, dok se ne isprave pogreske u tom poglavlju npr.:
%\includeonly{Poglavlje_1}


\begin{document}

\frontmatter   % - ne dirati

% upisati naziv studija
\degreesubject{Diplomski studij elektrotehnike}% upisati odgovarajuci naziv studija

% upisati vrstu rada
\documenttype{Diplomski rad}  % upisati ili Završni rad ili Diplomski rad

\title{Naslov rada}   % upisati specifični naslov rada

\date{\MONTH~\the\year.}   % ne dirati - mjesec i godina će se upisati sami

\author{Ime Prezime}  % upisati svoje ime i prezime
\jmbag{\hspace*{-4pt} 12345678901}  % upisati vlastiti JMBAG  (broj 2pt pod naredbom \hspace prilagoditi po potrebi da se JMBAG poravna s početkom imena - dodavanjem pozitivne ili negativne vrijednosti pt)
\maketitle		% ne dirati

%\makecopyright

% Okruzenje za pisanje posvete. Maknuti komentare ukoliko se želi napisati posvetu.
%\begin{dedication}
%	Ovo je posveta nekome
%\end{dedication}

\mentor{prof.~dr.~sc.~Miroslav Joler}   % zamijeniti podacima o svojem mentoru
\maketitleabstract

% kreira mjesto za umetnuti stranicu s opisom zadatka - ne dirati
\begin{assignmentpage}
\end{assignmentpage}

% kreira mjesto za umetnuti stranicu s izjavom o samostalnoj izradbi zadatka - ne dirati
\begin{honestystatementpage}
    \input{Izjava}
\end{honestystatementpage}

% Okruzenje za pisanje zahvale
\begin{acknowledgments} % staviti znak komentara ukoliko se ne stavlja tekst zahvale
    \input{Zahvala}  
\end{acknowledgments}

% kreiranje popisa sadrzaja, slika i tabela - nista ne dirati
\renewcommand{\baselinestretch}{1}\normalsize
\tableofcontents
\listoffigures
\listoftables
\renewcommand{\baselinestretch}{1.2}\normalsize

\mainmatter		% ne dirati

% Ovdje pomoću include funkcije ucitavati kreirana poglavlja. Poglavljima dajte logicna imena s obzirom na sadrzaj prikazan u njima (bez razmaka u imenu).

\include{Intro}   % ovo poslije staviti pod komentar kada se nauči koristiti
%\include{Poglavlje_1}  % dati neko logicno ime umjesto ``Poglavlje_1''
%\include{Poglavlje_2} 
%\include{Poglavlje_3} % itd.

% kreiranje popisa literature
%\bibliographystyle{tex_aux/IEEEtranHR}  % ne dirati
%% \bibliographystyle{ieeetr}
%% \bibliography{Literatura}  % ime .bib datoteke gdje vam je opisana literatura pomoću BiBTeX sintakse (npr. pomoću JabRef programa)

\include{Literatura}
%%%%%%%%%%%%%%%%%%%%%%%%%%%%%%%%%%%%%%%%%%%%%%%%%%%%%%
% kreiranje popisa literature koristeći BiBLaTeX
%\printbibliography[heading=bibintoc,title={Literatura}] % included in Literatura.tex
%%%%%%%%%%%%%%%%%%%%%%%%%%%%%%%%%%%%%%%%%%%%%%%%%%%%%%

%%%%%%%%%%%%%%%%%%%%%%%%%%%%%%%%%%%%%%%%%%%%%%%%%%%%%%
%:::::::::: ukljucenje popisa kratica u tekst ::::::::
% Blok linija koda ispod ovoga generira ukljucenje popisa kratica u tekstu. Za uporabu, vidjeti Upute.
% Nije obavezno. Ako se ne zeli koristiti, onda ovaj blok staviti u komentar pomocu znaka %
\input{Kratice}
\printglossary[type=\acronymtype, title=Pojmovnik]
%\pagestyle{plain}
%%%%%%%%%%%%%%%%%%%%%%%%%%%%%%%%%%%%%%%%%%%%%%%%%%%%%

%:::::::::::: blok za definiranje Sazetka/Abstracta rada 
\begin{abstract}
    \input{Sazetak}	  % sazetak rada i kljucne rijeci na HR i EN
\end{abstract}


%:::::::::::: PRILOZI (neobavezno) ::::::::::::::::::::
% ispod \appendix zaglavlja pomocu \include dodati poglavlja s prilozima
% ukoliko nemate priloga, ovaj blok linija staviti u komentar
\appendix
\include{Prilog_1}  % dati neko suvislo ime umjesto ovoga
%\include{Prilog_2}  % itd.


\end{document}


\end{document}